\section{Conclusion}
\label{sec:conclusion}

In this work, we investigate when teacher intervention improves the trajectories used for
on-policy distillation of multi-turn agents. Our controlled experiments show that a teacher
correction at a low-confidence step improves subsequent student behavior and task success,
outperforming random intervention. Building on this
finding, we propose UOPD, which uses teacher uncertainty to select turns for action correction
and imitation learning while retaining standard OPD on student actions. Experiments on
ALFWorld, WebShop, and multi-hop Search demonstrate that UOPD improves task
performance over OPD while reducing interaction turns.

Our findings highlight the value of teacher uncertainty for guiding both what an agent learns
and the trajectories from which it learns. Selective corrections provide useful learning targets
and steer subsequent interaction toward better outcomes. Future work could extend this approach
to longer-horizon tasks and adapt the intervention budget to task difficulty and student progress.
